%% ─────────────────────────────────────────────────────────────────────────────
%% Lettre de motivation — Talvin Ackbaraly × MBDA France
%% Stage : apport du deep learning pour la segmentation de cibles aériennes
%% Même charte que le CV (XCharter, un seul accent, filets fins)
%% Compile: latexmk -pdf lettre-motivation-mbda.tex
%% ─────────────────────────────────────────────────────────────────────────────
\documentclass[11pt, a4paper]{extarticle}

\usepackage[T1]{fontenc}
\usepackage[utf8]{inputenc}
\usepackage[french]{babel}
\usepackage[top=1.6cm, bottom=1.6cm, left=2.2cm, right=2.2cm]{geometry}
\usepackage{XCharter}
\usepackage[letterspace=80]{microtype}
\usepackage{ragged2e}
\usepackage{xcolor}
\usepackage[hidelinks]{hyperref}

\hypersetup{pdftitle={Lettre de motivation — Talvin Ackbaraly}, pdfauthor={Talvin Ackbaraly}}

%% ── Colour system: identique au CV ──────────────────────────────────────────
\definecolor{accent}{HTML}{2563A0}
\definecolor{ink}{HTML}{1A1A1A}
\definecolor{soft}{HTML}{6E6E6E}
\definecolor{hair}{HTML}{C9C9C9}

\hypersetup{colorlinks=true, urlcolor=accent, linkcolor=accent}
\AtBeginDocument{\color{ink}}

\pagestyle{empty}
\setlength{\parindent}{0pt}
\setlength{\JustifyingParindent}{0pt}
\setlength{\parskip}{8pt}
\linespread{1.08}
\microtypesetup{expansion=false}

\begin{document}

%% ── En-tête : expéditeur ─────────────────────────────────────────────────────
%% Coordonnées seulement : l'école et les liens sont déjà dans le CV
{\color{accent}\fontsize{22}{26}\selectfont\scshape\textls[40]{Talvin Ackbaraly}}\par\vspace{3pt}
{\setlength{\parskip}{0pt}
94270 Le Kremlin-Bicêtre\\
\href{tel:+33769389008}{\textcolor{ink}{07 69 38 90 08}}\\
\href{mailto:talvin.ackbaraly@icloud.com}{\textcolor{ink}{talvin.ackbaraly@icloud.com}}\par}
\vspace{-4pt}{\color{hair}\rule{\textwidth}{0.4pt}}\par

%% ── Destinataire et date ─────────────────────────────────────────────────────
\vspace{2pt}
%% Bloc calé contre la marge droite, texte aligné à gauche dans le bloc
\hfill\begin{tabular}[t]{@{}l@{}}
  \textbf{MBDA France}\\
  Direction Technique\\
  1 avenue Réaumur\\
  92350 Le Plessis-Robinson
\end{tabular}

Le \today

\vspace{2pt}
\textbf{Objet : candidature au stage de fin d'études « Apport du deep learning pour la segmentation
de cibles aériennes » — 6 mois à partir de février 2027}

\vspace{2pt}
Madame, Monsieur,


\justifying
Étudiant en dernière année du cycle ingénieur à l'EPITA, dans la majeure Image, je vous adresse
ma candidature pour le stage consacré à l'apport du deep learning pour la segmentation de cibles
aériennes, au sein de la Direction Technique de MBDA France. La segmentation est le domaine de la
vision que j'ai le plus pratiqué. Un sujet d'étude, où il faut comparer des méthodes, les mesurer
et justifier ses choix, correspond à la façon dont j'aime travailler.

J'ai implémenté un U-Net de A à Z en PyTorch, sans modèle pré-entraîné, pour segmenter premier
plan et arrière-plan sur COCO : architecture encodeur-décodeur, boucle d'entraînement et
évaluation sur des données non vues. Ce projet m'a appris ce que chaque choix d'architecture
change réellement. J'ai aussi travaillé sur la segmentation d'images médicales, des poumons sur
scanner CT et d'un gliome sur IRM, où la cible occupe souvent une petite partie de l'image. Je
pense retrouver cette difficulté avec des cibles aériennes, qui ne couvrent parfois que quelques
pixels sur un fond de ciel.

Mon projet de fin d'études, mené avec la Bibliothèque nationale de France, suit la démarche que
décrit votre offre. J'y suis en charge de la détection : j'ai comparé plusieurs architectures,
réglé les hyperparamètres et entraîné les modèles sur 10\,000 images annotées. Sur le jeu de
validation, j'obtiens 94,7\,\% de mAP50, mais 85,2\,\% en mAP50-95 : l'écart montre à quel
point le choix de la métrique change la lecture des résultats. Pour des cibles aériennes, un IoU
global ne suffirait pas, car une cible de quelques pixels y pèse très peu. J'aimerais donc
compléter les mesures classiques (IoU, Dice) par une analyse selon la taille des cibles et
la qualité des contours.

L'optimisation de la complexité du modèle m'intéresse aussi beaucoup. J'ai porté un traitement
vidéo du CPU au GPU en CUDA, avec une accélération de ×24, en analysant les goulots
d'étranglement avec Nsight. ONNX et TensorRT sont les outils que je souhaite maîtriser ensuite,
et ce stage serait l'occasion de les appliquer à un cas réel. Enfin, je suis à l'aise pour
rédiger un rapport scientifique, en français comme en anglais (TOEIC 845), ce qui compte dans
l'environnement international de MBDA.

Je suis disponible pour un stage de six mois à partir de février 2027 et serais heureux de vous
présenter plus en détail ma motivation lors d'un entretien.

Je vous prie d'agréer, Madame, Monsieur, l'expression de mes salutations distinguées.

\vspace{10pt}
\hfill Talvin Ackbaraly\hspace*{1cm}

\end{document}
